\chapter[Lezione 23 --- 21 novembre 2026]{Lezione 23: teorema di Bernoulli--de l'Hôpital e applicazioni}
\label{chap:lezione-23}

\section{Teorema di Bernoulli--de l'Hôpital}
\label{sec:lez23-hopital}

Il teorema di Bernoulli--de l'Hôpital permette di ricondurre
il limite di un rapporto indeterminato al limite del rapporto
tra le derivate, purché siano verificate le sue ipotesi.

\begin{theorem}[Teorema di Bernoulli--de l'Hôpital: forma \(0/0\)]
\label{thm:lez23-hopital-zero}
Siano \(a,b\in\R\), con \(a<b\), e siano
\(f,g\colon(a,b)\longrightarrow\R\) derivabili su \((a,b)\),
con \(g'(x)\ne0\) per ogni \(x\in(a,b)\).
Si supponga che
\[
  \lim_{x\to a^+}f(x)=\lim_{x\to a^+}g(x)=0
\]
e che esista finito il limite
\[
  \lim_{x\to a^+}\frac{f'(x)}{g'(x)}=\ell\in\R.
\]
Allora
\[
  \lim_{x\to a^+}\frac{f(x)}{g(x)}=\ell.
\]
\end{theorem}

L'enunciato vale anche per \(x\to b^-\) e per i limiti
unilateri in un punto interno al dominio, assumendo
le stesse ipotesi nell'intervallo da cui si tende al punto.
Se le ipotesi valgono da entrambi i lati e il rapporto
tra le derivate ammette lo stesso limite, il risultato
vale anche per il limite bilatero.
La dimostrazione del teorema non viene sviluppata in questa sede.

\begin{remark}[Condizioni di applicabilità]
\label{rem:lez23-hopital-condizioni}
Prima di applicare il \cref{thm:lez23-hopital-zero}
è necessario verificare che il rapporto sia della forma
indeterminata \(0/0\), che il denominatore derivato non
si annulli e che esista il limite del rapporto tra le derivate.
Se quest'ultimo limite non esiste, il teorema non permette
di concludere che non esista neppure il limite iniziale.
Inoltre non si può usare il teorema per dimostrare un limite
da cui dipende la formula di derivazione impiegata:
per esempio, non è lecito dedurre
\(\lim_{x\to0}\sin x/x=1\) utilizzando
\((\sin x)'=\cos x\), poiché tale formula è stata
ricavata proprio da quel limite nel
\cref{thm:lez20-limite-seno}.
\end{remark}

\section{Esempi nella forma \(0/0\)}
\label{sec:lez23-esempi-zero}

\begin{example}[Limite di \((x-\sin x)/x^3\)]
\label{ex:lez23-seno-cubico}
Si consideri
\[
  \lim_{x\to0^+}\frac{x-\sin x}{x^3}.
\]
Il numeratore e il denominatore tendono a \(0\).
Per \(x>0\), posto \(f(x):=x-\sin x\) e
\(g(x):=x^3\), si ha
\[
  f'(x)=1-\cos x,
  \qquad g'(x)=3x^2\ne0.
\]
Per la \cref{prop:lez20-limite-coseno},
\[
\begin{aligned}
  \lim_{x\to0^+}\frac{f'(x)}{g'(x)}
  &=\frac13\lim_{x\to0^+}\frac{1-\cos x}{x^2}\\
  &=\frac13\cdot\frac12=\frac16.
\end{aligned}
\]
Pertanto, per il \cref{thm:lez23-hopital-zero},
\[
  \lim_{x\to0^+}\frac{x-\sin x}{x^3}=\frac16.
\]
\end{example}

\begin{example}[Il limite del rapporto può esistere senza il limite delle derivate]
\label{ex:lez23-hopital-non-applicabile}
Si consideri
\[
  \lim_{x\to0^+}\frac{x^2\sin(1/x)}{2x+\sin x}.
\]
Posti \(f(x):=x^2\sin(1/x)\) e
\(g(x):=2x+\sin x\), entrambi i termini tendono
a \(0\) per \(x\to0^+\). Inoltre, per \(x>0\),
\[
  f'(x)=2x\sin(1/x)-\cos(1/x),
  \qquad g'(x)=2+\cos x\ge1.
\]
Tuttavia il limite di \(f'(x)/g'(x)\) non esiste.
Infatti, scegliendo
\[
  a_n:=\frac1{2\pi n},
  \qquad b_n:=\frac1{(2n+1)\pi},
  \qquad n\ge1,
\]
si ottiene
\[
  \lim_{n\to+\infty}\frac{f'(a_n)}{g'(a_n)}=-\frac13,
  \qquad
  \lim_{n\to+\infty}\frac{f'(b_n)}{g'(b_n)}=\frac13.
\]
Il teorema non è quindi applicabile. Il limite richiesto,
però, esiste: per \(x>0\) si può scrivere
\[
  \frac{x^2\sin(1/x)}{2x+\sin x}
  =\frac{x\sin(1/x)}{2+\dfrac{\sin x}{x}}.
\]
Poiché \(|x\sin(1/x)|\le x\to0\) e,
per il \cref{thm:lez20-limite-seno},
\(2+\sin x/x\to3\), si conclude
\[
  \lim_{x\to0^+}\frac{x^2\sin(1/x)}{2x+\sin x}=0.
\]
\end{example}

\begin{example}[Limite esponenziale con parametro]
\label{ex:lez23-esponenziale-parametro}
Sia \(\alpha>0\). Si determini
\[
  \lim_{x\to0^+}
  \frac{e^x-1-x-\frac{x^2}{2}}{x^\alpha}.
\]
Si consideri dapprima il caso \(\alpha=3\).
Applicando tre volte il \cref{thm:lez23-hopital-zero}
si ottiene
\[
\begin{aligned}
  \lim_{x\to0^+}\frac{e^x-1-x-\frac{x^2}{2}}{x^3}
  &=\lim_{x\to0^+}\frac{e^x-1-x}{3x^2}\\
  &=\lim_{x\to0^+}\frac{e^x-1}{6x}\\
  &=\lim_{x\to0^+}\frac{e^x}{6}=\frac16.
\end{aligned}
\]
Posto
\[
  A(x):=\frac{e^x-1-x-\frac{x^2}{2}}{x^3},
\]
si ha quindi \(A(x)\to1/6>0\) e
\[
  \frac{e^x-1-x-\frac{x^2}{2}}{x^\alpha}
  =A(x)x^{3-\alpha}.
\]
Distinguendo i casi,
\[
  \lim_{x\to0^+}
  \frac{e^x-1-x-\frac{x^2}{2}}{x^\alpha}
  =\begin{cases}
    0, & 0<\alpha<3,\\[2pt]
    \dfrac16, & \alpha=3,\\[4pt]
    +\infty, & \alpha>3.
  \end{cases}
\]
\end{example}

\begin{corollary}[Criterio di derivabilità mediante la derivata]
\label{cor:lez23-limite-derivata}
Sia \(f\colon I\longrightarrow\R\), con \(I\)
intervallo e \(x_0\) interno a \(I\). Si supponga
che \(f\) sia continua in \(x_0\), derivabile
in un intorno bucato di \(x_0\) e che
\[
  \lim_{x\to x_0}f'(x)=\ell\in\R.
\]
Allora \(f\) è derivabile in \(x_0\) e
\[
  f'(x_0)=\ell.
\]
La conclusione vale anche per una sola derivata
unilaterale, assumendo le ipotesi sul lato corrispondente.
\end{corollary}

\begin{proof}
Per la continuità di \(f\) in \(x_0\),
\(f(x)-f(x_0)\to0\) per \(x\to x_0\).
Applicando il \cref{thm:lez23-hopital-zero}
separatamente ai due limiti unilateri, con
\(g(x):=x-x_0\) e \(g'(x)=1\), si ha
\[
  \lim_{x\to x_0^\pm}
  \frac{f(x)-f(x_0)}{x-x_0}
  =\lim_{x\to x_0^\pm}f'(x)=\ell.
\]
I due limiti coincidono; per la
\cref{def:lez22-derivabilita} segue la tesi.
\end{proof}

\section{Forme del teorema agli estremi e all'infinito}
\label{sec:lez23-hopital-infinito}

Il teorema ammette versioni analoghe per i limiti
all'infinito e per la forma indeterminata
\(\infty/\infty\).

\begin{theorem}[Forma \(0/0\) per \(x\to+\infty\)]
\label{thm:lez23-hopital-zero-infinito}
Siano \(a\in\R\) e \(f,g\colon(a,+\infty)\longrightarrow\R\)
derivabili, con \(g'(x)\ne0\) per ogni \(x>a\).
Se
\[
  \lim_{x\to+\infty}f(x)
  =\lim_{x\to+\infty}g(x)=0,
  \qquad
  \lim_{x\to+\infty}\frac{f'(x)}{g'(x)}=\ell\in\R,
\]
allora
\[
  \lim_{x\to+\infty}\frac{f(x)}{g(x)}=\ell.
\]
\end{theorem}

\clearpage

\begin{theorem}[Forma \(\infty/\infty\) a un estremo finito]
\label{thm:lez23-hopital-infinito-estremo}
Siano \(a,b\in\R\), con \(a<b\), e siano
\(f,g\colon(a,b)\longrightarrow\R\) derivabili,
con \(g'(x)\ne0\) per ogni \(x\in(a,b)\).
Se
\[
  \lim_{x\to a^+}f(x)
  =\lim_{x\to a^+}g(x)=+\infty,
  \qquad
  \lim_{x\to a^+}\frac{f'(x)}{g'(x)}=\ell\in\R,
\]
allora
\[
  \lim_{x\to a^+}\frac{f(x)}{g(x)}=\ell.
\]
\end{theorem}

\begin{theorem}[Forma \(\infty/\infty\) per \(x\to+\infty\)]
\label{thm:lez23-hopital-infinito}
Siano \(a\in\R\) e \(f,g\colon(a,+\infty)\longrightarrow\R\)
derivabili, con \(g'(x)\ne0\) per ogni \(x>a\).
Se
\[
  \lim_{x\to+\infty}f(x)
  =\lim_{x\to+\infty}g(x)=+\infty,
  \qquad
  \lim_{x\to+\infty}\frac{f'(x)}{g'(x)}=\ell\in\R,
\]
allora
\[
  \lim_{x\to+\infty}\frac{f(x)}{g(x)}=\ell.
\]
\end{theorem}

Per le forme precedenti valgono versioni analoghe
per \(x\to-\infty\), per \(x\to b^-\) e, nel caso
di un estremo finito, per limiti unilateri in un
punto interno. Le dimostrazioni non vengono
sviluppate in questa sede.

\section{Confronti di crescita}
\label{sec:lez23-confronti-crescita}

\begin{proposition}[Il logaritmo cresce meno di ogni potenza]
\label{prop:lez23-logaritmo-potenza}
Per ogni \(\varepsilon>0\),
\[
  \lim_{x\to0^+}x^\varepsilon\log x=0.
\]
\end{proposition}

\begin{proof}
Poiché \(-\log x\to+\infty\) e
\(x^{-\varepsilon}\to+\infty\), si scriva
\[
  x^\varepsilon\log x
  =-\frac{-\log x}{x^{-\varepsilon}}.
\]
Per \(x>0\), il denominatore derivato è
\(-\varepsilon x^{-\varepsilon-1}\ne0\).
Dal \cref{thm:lez23-hopital-infinito-estremo} segue
\[
\begin{aligned}
  \lim_{x\to0^+}\frac{-\log x}{x^{-\varepsilon}}
  &=\lim_{x\to0^+}
    \frac{-1/x}{-\varepsilon x^{-\varepsilon-1}}\\
  &=\frac1\varepsilon\lim_{x\to0^+}x^\varepsilon=0.
\end{aligned}
\]
Pertanto anche \(x^\varepsilon\log x\to0\).
\end{proof}

\begin{proposition}[L'esponenziale cresce più di ogni potenza]
\label{prop:lez23-esponenziale-potenza}
Per ogni \(\alpha>0\),
\[
  \lim_{x\to+\infty}\frac{e^x}{x^\alpha}=+\infty.
\]
\end{proposition}

\begin{proof}
Si scelga \(k:=\lfloor\alpha\rfloor+1\in\N\),
così che \(k>\alpha\).

Per \(k\ge2\), le derivate successive di \(x^k\)
sono
\[
\begin{aligned}
  (x^k)' &= kx^{k-1},\\
  (x^k)'' &= k(k-1)x^{k-2},\\
  &\ \vdots\\
  (x^k)^{(k)} &= k(k-1)\cdots2\cdot1=k!.
\end{aligned}
\]
La formula finale vale anche per \(k=1\).

Poiché ogni derivata di \(e^x\) coincide con
\(e^x\), applicando \(k\) volte il
\cref{thm:lez23-hopital-infinito} si ottiene
\[
\begin{aligned}
  \lim_{x\to+\infty}\frac{x^k}{e^x}
  &=\lim_{x\to+\infty}\frac{kx^{k-1}}{e^x}\\
  &=\cdots
  =\lim_{x\to+\infty}\frac{k!}{e^x}=0.
\end{aligned}
\]
Ne segue, essendo i rapporti positivi,
\[
  \lim_{x\to+\infty}\frac{e^x}{x^k}=+\infty.
\]
Infine, poiché \(k-\alpha>0\),
\[
  \frac{e^x}{x^\alpha}
  =\frac{e^x}{x^k}\,x^{k-\alpha}
  \longrightarrow+\infty.
\]
\end{proof}

\section{Esercizio su continuità e derivabilità}
\label{sec:lez23-esercizio-tratti}

\begin{exercise}[Funzione definita a tratti con tre parametri]
\label{ex:lez23-continuita-derivabilita-tratti}
Siano \(a,b,c\in\R\), con \(b\ne0\) e \(c\ge0\).
Si consideri la funzione
\[
  f(x):=
  \begin{cases}
    \dfrac{\log(x-1)}{b(x-2)}, & x>2,\\[6pt]
    a, & x=2,\\[4pt]
    \arctan\!\left(\dfrac{c}{x-2}\right), & x<2.
  \end{cases}
\]
Determinare per quali valori dei parametri
\(f\) è continua e derivabile su \(\R\).
\end{exercise}

\begin{proof}[Soluzione dell'esercizio]
La funzione è continua nei due intervalli
\((-\infty,2)\) e \((2,+\infty)\).
Resta da studiare il comportamento in \(x_0=2\).

\emph{Primo punto: continuità.}
Per il limite notevole del logaritmo,
posto \(t:=x-2\), si ha
\[
\begin{aligned}
  \lim_{x\to2^+}f(x)
  &=\frac1b\lim_{t\to0^+}\frac{\log(1+t)}t
   =\frac1b.
\end{aligned}
\]
Dal lato sinistro, invece,
\[
  \lim_{x\to2^-}f(x)
  =\begin{cases}
    0, & c=0,\\[2pt]
    -\dfrac\pi2, & c>0.
  \end{cases}
\]
Poiché \(b\ne0\), il limite destro \(1/b\)
non può essere nullo. Pertanto la continuità
in \(2\) è possibile soltanto quando \(c>0\)
e
\[
  \frac1b=a=-\frac\pi2.
\]
Si conclude che
\[
  f\text{ è continua su }\R
  \quad\Longleftrightarrow\quad
  a=-\frac\pi2,\qquad b=-\frac2\pi,\qquad c>0.
\]

\emph{Secondo punto: derivabilità.}
Per il \cref{thm:lez22-derivabilita-continuita},
la continuità è necessaria alla derivabilità.
Si assumano dunque \(a=-\pi/2\), \(b=-2/\pi\)
e \(c>0\), e si calcolino le derivate unilaterali
in \(2\).

Da destra, usando ancora \(t:=x-2\),
\[
\begin{aligned}
  \lim_{x\to2^+}\frac{f(x)-f(2)}{x-2}
  &=\frac\pi2\lim_{t\to0^+}
    \frac{t-\log(1+t)}{t^2}\\
  &=\frac\pi2\lim_{t\to0^+}
    \frac{1-\frac1{1+t}}{2t}\\
  &=\frac\pi4\lim_{t\to0^+}\frac1{1+t}
   =\frac\pi4.
\end{aligned}
\]
Nel secondo passaggio si è applicato il
\cref{thm:lez23-hopital-zero}.

Da sinistra, per \(c>0\), si ha una forma \(0/0\).
Poiché
\[
  \frac{d}{dx}\arctan\!\left(\frac{c}{x-2}\right)
  =-\frac{c}{(x-2)^2+c^2},
\]
dal \cref{thm:lez23-hopital-zero} segue
\[
\begin{aligned}
  \lim_{x\to2^-}\frac{f(x)-f(2)}{x-2}
  &=\lim_{x\to2^-}
    \frac{\arctan\!\left(\frac{c}{x-2}\right)+\frac\pi2}
         {x-2}\\
  &=\lim_{x\to2^-}\frac{-c}{(x-2)^2+c^2}
   =-\frac1c<0.
\end{aligned}
\]
Le due derivate unilaterali non coincidono, poiché
\(\pi/4>0\) e \(-1/c<0\). Pertanto \(f\)
non è derivabile in \(2\) per nessun valore
ammissibile dei parametri e, di conseguenza,
non è mai derivabile su tutto \(\R\).
\end{proof}
